
\kospis{Ciąg Marka--Ulama}
{Bartłomiej Pawlik} 

\vspace*{-24pt}

\kotyt4{Ciąg Marka--Ulama}


\vspace*{-4pt}
\marg[20]{
  \waut{Bartłomiej PAWLIK} %\hfill\scriptsize\rm

  \medskip
    Politechnika Śląska}
%\begin{dwieszpalty}
Potrzeba szybkiego przekazywania współpracownikom świeżych i~istotnych
informacji jest oczywiście starsza niż internet i~e-maile. W~latach 40. XX wieku
w~laboratorium w~Los Alamos\marg{W~czasie II wojny światowej w~Los Alamos
znajdowała się siedziba Projektu Manhattan, który miał na celu stworzenie
pierwszej na świecie bomby atomowej. Cel ten został osiągnięty w~roku 1945, a~opracowana broń
została użyta w~ataku na Hiroszimę dokładnie 21 dni po pierwszej udanej
próbie przeprowadzonej na pustyni w~Nowym Meksyku. Po wojnie placówka została
odtajniona i~w~dalszym ciągu prowadzono tam prace, m.in. nad konstrukcją bomby
wodorowej.} radzono sobie z~nią poprzez rozpowszechnianie kopii krótkich
notatek. Dzięki temu pracownicy naukowi placówki mogli być na bieżąco z~wynikami
swoich kolegów. Niemałą konsternację i~rozbawienie wywołała (zamieszczona
niżej) notatka
rozesłana pewnego grudniowego dnia 1947 roku.\vspace*{-6pt}
$$
\includegraphics[scale=.35]{memo.png}
$$
Notatka zawiera drobną łamigłówkę. Przedstawiona w~niej sekwencja to tak naprawdę ciąg liczb od 0 do 100 ustawionych w~porządku alfabetycznym względem ich nazw w~języku angielskim. Rozszyfrowanie tego wzorca utrudniał fakt, że ciąg w~sporządzonej w~pośpiechu notatce zawierał błędy: Brakuje w~nim liczby 10, natomiast liczba 12 występuje dwa razy (jako $\mathtt{dozen}$ i~jako $\mathtt{twelve}$).

Jeden z~członków dyrekcji stwierdził,\marg{Przedstawione informacje  pochodzą  z~notatki ,,A~Memorable Memo''
zamieszczonej w~\textit{Los Alamos Science} Nr~15, 1987.}
że autora notatki należy zwolnić za niepoważne zachowanie. Więcej zrozumienia
wykazał rozbawiony przedstawiciel waszyngtońskiej Komisji Energii Atomowej,
który stwierdził, że tekścik to {\it najlepsza rzecz, jaka dotychczas została
stworzona w~Los Alamos}.
Figlarzami odpowiedzialnymi za sformułowanie notatki okazali się J. Carson Mark i~Stanisław Ulam.
\vspace*{\parskip}
\begin{dwieszpalty}
Czytelnika Zaciekawionego zachęcamy do rozwiązania dwóch prostych zadań
(odpowiedzi na s.~\hypertarget{pytania}{}\hyperlink{odp}{9}):
\begin{enumerate}[label=(\alph*)]
\item Między które dwie liczby w~powyższym ciągu należy wstawić brakującą 10?
\item W~polskojęzycznej wersji ciągu Marka--Ulama ostatnim elementem jest -- tak
jak w~oryginale -- liczba 0. Jakie liczby są, odpowiednio, pierwszym i~przedostatnim elementem tego ciągu?	
\end{enumerate}

Na koniec dodajmy, że jest to jeden z~nielicznych przykładów ciągu figurującego w~OEIS, który celowo zawiera błędy. Ciąg Marka--Ulama ma numer \href{https://oeis.org/A261402}{A261402}, a~jego poprawiona wersja to \href{https://oeis.org/A036747}{A036747}.
\end{dwieszpalty}
\endinput	


%\end{dwieszpalty}





